\part{Unchecked targets and remaining gates}
\label{part:open}
\section{What has not been proved}
There is no unchecked finding offered as a theorem in this part. The
following are the unresolved targets of the task. Their state is a
current gate, not a proved negative answer. A precise counterexample to
one proposed transfer is recorded separately from a proof that no other
transfer can exist.

\begin{description}
\item[O-PH3-WEAK: numerical positive-form clause for irreducible-inner plethysm.]
P-F19 refutes the stronger canonical modular-index clause of
Hypothesis 1.6.7, already with inner shape $(3)$: every h.s.o.p.
in the displayed family has an exponential degree. The first,
weaker clause asks only for some positive rational expression of
polynomial modular index. It does not require that expression to
come from parameters of the canonical ring, and is not refuted by
P-F19. P-F11--P-F12 show that this weaker bound would suffice for
the revised PH2. P-F14--P-F15 separately refute both index bounds
in the broad rank-only subgroup-restriction formulation. The
remaining task is to prove or refute the numerical positive-form
bound, PH2 and the narrower saturation bounds, with their exact
family and parameter restrictions.

P-F20 proves an exponential universal factor and an actual pole in
some fixed-target covariant module, but neither conclusion by itself
gives an ordinary stretching pole. P-F21--P-F22 give exact
counterexamples to general quantitative parameter-to-Hilbert
transfers; P-F20 also defeats an exact-factor transfer within
factorial rational invariant rings of a connected semisimple group.
The stronger surviving geometric route is now explicit. By P-F24
the chain gives a smooth one-dimensional component of the grading
$\boldsymbol\mu_{2^{n-1}}$ fixed scheme. By P-F23 an actual
primitive-root pole follows if this is its unique maximal-dimensional
component. It remains to exclude or control competing components
represented by singular polystable cubics. P-F26 proves the existence
of a finite two-power lift and confines all competitors to the
discriminant boundary. P-F27 handles arbitrary lift order for graph
supports and forces every singular competitor to contain a monomial
in three distinct variables in every lift eigenbasis. The remaining
support problem is therefore a cubic hypergraph problem. The
connected-graph extension of P-F27 further forces its graph part
to be disconnected, and P-F29 forces both square-linear and
three-distinct-variable monomials. P-F28 computes the chain
coefficient and its two-adic valuation exactly, without asserting
a comparison with singular contributions. A phase
calculation proving noncancellation without uniqueness would also suffice.

\item[O-COUNT: uniform effective shifted counting.]
P-F2 constructs a single shifted polytope from explicit positive Hilbert
data. P-F12 controls their existence and size under the weaker
positive-form hypothesis in expanded parameters. P-F13 reconstructs
and verifies candidate data from a finite prefix under an explicit
denominator promise. A procedure constructing them with the original
partition-input complexity bounds remains to be given. A small object
whose existence is known separately for every instance is not automatically
a polynomial-time constructible object. The homogeneous version is already
refuted by P-K1; no revision should silently restore that false assertion.

\item[O-NSQ: nonstandard geometric realization.]
For the particular bases and products in GCT VII--VIII, one must construct
the category, identify the basis, control parity or shifts, establish
product and specialization compatibility, and supply size bounds.
P-F6 is a sufficient criterion after these data exist. The determinant
and theta-lift passages of the supplied note do not themselves identify
the relevant Hom spaces. P-F5 refutes the simplest attempted substitution
of ordinary interlacing for plethysm restriction.

\item[O-OBSTRUCTION: an actual family of separating equations.]
In the no-occurrence range, prove a strict inequality between the actual
coordinate-ring multiplicities, or construct an explicit polynomial
vanishing on the determinant closure and not on the padded permanent.
The ambient plethysm and Kronecker multiplicities do not determine the
ranks of the quotient maps. P-F8 and PL-06 isolate this rank/comparison
requirement. No such asymptotic obstruction family is constructed here.

\item[O-DIR-HN: arbitrary solution fields in Nullstellensatz.]
C-GRH-01 proves only the cyclotomic-field transfer. Koiran's algorithm uses
a uniform prime-ideal/root-modulo-prime estimate for arbitrary primitive
solution fields. Rojas's DZH or GIPIT conditions are alternative sufficient
analytic inputs. The present supplied Dirichlet theorem does not state
these inputs. No formal nonimplication between the analytic conjectures
is proved.

\item[O-PIT: the derandomization required by GCT V.]
The exact algorithm or hardness hypothesis of each theorem must still be
provided. The supplied counting algorithm yields C-PIT-01 in a bounded
finite-field circuit model. It does not supply the characteristic-zero
black-box, border or strengthened symbolic determinant PIT assertions
used elsewhere. GRH does not replace those separate hypotheses.

\item[O-INTEGRAL: uniform bounds for specialization.]
B7.1 proves good-prime survival for a fixed integral separator; B7.2 gives
an exact rank test in any fixed degree and integer minors excluding bad
primes. An efficient GCT transfer needs a uniform separator-degree bound,
succinct rank data and a usable prime bound. Neither Hilbert-scheme
representability nor the pointwise specialization theorem supplies these
quantitative statements.

\item[O-COMPLEXITY: Boolean and algebraic separation.]
The ETH theorem refutes the precise exponential hypotheses treated in
C-SAT-01. Its supplied circuit corollary is superlinear for
$\mathsf E^{\mathsf{NP}}$, not superpolynomial for $\mathsf{NP}$.
The current constant-free integer Boolean--arithmetic disjunction cited
in the complexity section does not choose a branch and does not identify
the usual complex-field $\mathsf{VP}$ with its constant-free integer
variant. The stated upper bounds and conditional bridges do not settle
$\mathsf P$ versus $\mathsf{NP}$ or the target GCT obstruction.
\end{description}

\section{Approaches tried and the next discriminating steps}
The qualitative-sign route was tested by actual determinant projectors
and a virtual-character counterexample (P-F3--P-F4). Its next step is an
identity with the intended multiplicity projector together with a size
bound. The standard-branching route was tested on $\operatorname{Sym}^2$
(P-F5); a replacement must use the actual Schur-functor homomorphism.

The purity route was tested at three increasingly specific levels:
semisimple categorical structure (P-F6), pure finite schemes encoding
\#3SAT (P-F7), and the actual canonical plethysm algebra carrying the
obstructing Hilbert function, a positive multiplicative basis and pure
intersection cohomology (PLE-F7--PLE-F8). The next step is an
equivariant graded comparison theorem, not another proof that an unrelated
cohomology group is pure. P-F10 separately tests the missing quantitative
control of a grading.

The Hilbert-series route made positive progress: P-F1 constructs the
optimal shifted model on the obstructing ray; PLE-F5 realizes its normal
invariant ring; PLE-F6 classifies all rescalings that repair homogeneous
Ehrhart realization; P-F2 treats every explicit positive rational form;
P-F9 and P-F11 supply coefficient-positivity bounds; P-F12 connects the
last bound to the numerical positive-form clause of revised PH3.
P-F13 gives finite reconstruction under a denominator promise;
P-F16--P-F18 give finite fixed-embedding and fixed-dimension bounds;
P-F19 refutes polynomial growth of the canonical parameter bound.
The next step is the weaker numerical positive-form question or a
direct PH2 argument, together with uniform construction. An h.s.o.p.
lower bound is not treated as a pole lower bound for the ordinary
Hilbert series.
P-F20 confirms that distinction with universal denominators and
fixed-target covariants. Its multivariate counterexample defeats a
transfer based only on quasi-polynomiality from an affine slice to
simultaneous rays. P-F21--P-F22 then defeat generic normal,
Cohen--Macaulay, rational and Gorenstein attempts to repair the
parameter-degree inference. P-F23--P-F24 make the next gate
strictly more concrete: the known chain component is smooth, and
the unresolved work is the contribution of the other fixed
components. P-F26 identifies the entire smooth fixed locus with
the chain orbit; P-F27 excludes every graph-support boundary
competitor, including rank-deficient supports and lifts of larger
two-power order. The binary determinant subfamily is separately treated
in P-F25: its realized holes are bounded by thirteen, but the
source's strict saturation definition also tests off-residue
constituent values and is not settled by that table.

The archive route distinguishes the three products on representations,
separates the two Burnside operations, checks the transcribed identities
on small groups, and treats the modular and quadratic-form caveats
(PL-01--PL-08). Its obstruction gate is the actual rank of the
Foulkes--Howe multiplicity-space map. A Schur-positive character of its
source and target does not compute that rank.

The norm and approximation routes were tested against scaling, exact
zero detection, adjacent integer multiplicities, and an explicit rank-two
family approaching rank one in normalized norm (B4.1--B5.3). They do give
valid contraction-model obstructions under the stated operator bounds.
A full algebraic-orbit obstruction requires an invariant and a certified
separating gap that survive the allowed degenerations.

The arithmetic and ETH routes track the exact analytic family, all input
and reduction sizes, the direction of reductions, arithmetic precision,
and the class in which a lower bound is obtained. Their next steps are
specified in O-DIR-HN, O-PIT and O-COMPLEXITY. These are not resolved by
relabeling the supplied subexponential bound as polynomial time.
